\section{Statement-to-declaration index}
\label{app:index}

This appendix gives repository-relative identifiers for the claims in the
main text.  Upstream denotes the pinned Cubical or separate HoTT Reals
checkout; generic CQTS denotes the pre-existing torsor development.

\begingroup\small\raggedright
\begin{description}[font=\normalfont\bfseries]
\item[HoTT Reals (upstream)]
\texttt{HoTTReals/Data/Real/Base.agda}: \agda{ℝ}.
\texttt{HoTTReals/Data/Real/Properties.agda}: \agda{ℝ-isSet}.
\texttt{HoTTReals/Data/Real/Algebra/Multiplication.agda}:
\agda{ℝCommutativeRing}.

\item[Constructive complex ring (new local)]
\texttt{src/ConstructiveComplex/Complex.agda}: \agda{Complex},
\agda{complexZero}, \agda{complexOne}, \agda{complexAdd},
\agda{complexNeg}, \agda{complexMul}, \agda{imaginaryUnit},
\agda{complexConjugate}, \agda{complexIsSet}, \agda{complexIsCommRing},
\agda{ComplexCommRing}.

\item[Algebraic units (new local)]
\texttt{src/ConstructiveComplex/Units.agda}: \agda{ComplexUnits},
\agda{complexUnitOne}, \agda{complexUnitMul}, \agda{complexUnitInv},
\agda{complexUnitsIsSet}, \agda{ComplexUnitsGroup},
\agda{complexUnitsRegularLeftTorsor}.

\item[Generic torsor construction (CQTS)]
\texttt{src/TorsorDelooping/Regular.agda}: \agda{regularLeftAction},
\agda{regularLeftTorsor}.  \texttt{src/TorsorDelooping/Moduli.agda}:
\agda{LeftTorsorModuli}, \agda{TorsorEquiv}, \agda{torsorPathEquiv},
\agda{fromLeftTorsor}.

\item[Generic loop recovery (CQTS)]
\texttt{src/TorsorDelooping/Delooping.agda}: \agda{torsorBase},
\agda{torsorLoopEquivAut}, \agda{regularAutToGroup},
\agda{rightRegularAut}, \agda{regularAutIsoGroup},
\agda{torsorLoopEquivGroup}.

\item[Constructive-complex delooping (new local)]
\texttt{src/ConstructiveComplex/Delooping.agda}: \agda{BComplexUnits},
\agda{bComplexUnits}, \agda{complexUnitsLoopEquiv}.

\item[Key Cubical interfaces (upstream)]
\texttt{Cubical/Algebra/CommRing/Base.agda}: \agda{makeIsCommRing}.
\texttt{Cubical/Tactics/CommRingSolver.agda}: \agda{solve!}.
\texttt{Cubical/Algebra/CommRing/Properties.agda}: \agda{Units}.
\texttt{Cubical/Algebra/Group/Base.agda}: \agda{makeGroup-right}.
\texttt{Cubical/Algebra/Group/Properties.agda}: \agda{·GroupAutomorphismR}.
\texttt{Cubical/Foundations/SIP.agda}: \agda{SIP}.
\end{description}
\endgroup

The mathematical explanations in the main text explain checked constructions;
they do not assert that every sentence is a separately named theorem.
